% Generated from docs/REPORT_DRAFT.md; edit the manuscript then rerun the converter.
\chapter{Thực nghiệm và kết quả}

\section{Nguyên tắc báo cáo}

Mọi kết quả phải gắn với backend, movie count, corpus type, K, iteration và cấu
hình. Silver corpus không gọi là gold/human-reviewed trước khi review gate qua.
QA và performance claims chỉ lấy từ production Neo4j artifacts.

\section{Dataset và graph}
\reportfigure{quality_metrics.pdf}{Tổng hợp chỉ số chất lượng dữ liệu và graph}{fig:quality-metrics}

Snapshot đầu vào có 2.001 record; pipeline loại minh bạch một Movie không có bất
kỳ quan hệ graph nào và giữ đúng 2.000 Movie hợp lệ. Exact IMDb ratings match
1.783/1.855 phim có IMDb ID. Graph sau reasoning có 37.349 node và 353.915
relationship, với zero
orphan, duplicate stable ID, missing required property và invalid relationship.

Các số liệu này lấy từ \texttt{data/processed/manifest.json} và
\texttt{experiments/results/neo4j\_validation.json} của cùng source checksum
\texttt{eeeb7e27...2c800}.

\begin{longtable}{|>{\raggedright\arraybackslash}p{0.44\textwidth}|>{\raggedright\arraybackslash}p{0.44\textwidth}|}
\caption{Tổng hợp dataset và graph}\label{tab:9-1} \\
\hline
\textbf{Chỉ số snapshot/graph} & \textbf{Giá trị cuối} \\ \hline
\endfirsthead\hline \textbf{Chỉ số snapshot/graph} & \textbf{Giá trị cuối} \\ \hline\endhead
Record TMDB đầu vào & 2.001 \\ \hline
Movie hợp lệ sau quality gate & 2.000 \\ \hline
Movie bị loại do không có quan hệ & 1 \\ \hline
Movie có IMDb ID & 1.855 (92,75\%) \\ \hline
IMDb rating exact-match & 1.783/1.855 (96,12\%) \\ \hline
Movie có cast & 99,00\% \\ \hline
Movie có director & 99,30\% \\ \hline
Movie có genre & 98,60\% \\ \hline
Orphan/duplicate/missing/invalid/unsupported fact sau import & 0 \\ \hline
\end{longtable}

Quality gate không xóa im lặng record lỗi: Movie bị loại được ghi trong
\texttt{invalid\_records} với \texttt{reason=no\_graph\_relationships} và source row/TMDB ID, nhờ
đó tổng đầu vào, tổng hợp lệ và lý do chênh lệch vẫn truy vết được.

\section{Entity resolution}

Corpus silver gồm 100 case: 75 positive và 25 negative, có source ID/evidence.
Kết quả là TP=75, TN=25, FP=0, FN=0, do đó precision, recall và F1 đều bằng 1,00.
Kết quả chứng minh workflow xử lý đúng corpus deterministic này; nó không chứng
minh mọi tên thực tế được giải quyết hoàn hảo. Muốn gọi human-reviewed, mỗi case
phải có reviewer độc lập, timestamp, decision, rubric và adjudication note.

\section{Reasoning}

Năm mươi fact \texttt{CO\_STARRED\_WITH} silver được kiểm tra bằng supporting movie và
source cast; 50/50 fact hợp lệ, precision bằng 1,00. Validation bổ sung còn kiểm
tra mọi derived edge phải có ít nhất một shared Movie. Semantic workflow trên
full normalized snapshot materialize 35.419 triple, tăng từ 154.970 lên 190.389
triple và không có semantic violation. RDF export dùng chính processed snapshot
2.000 Movie như Neo4j; cả 10 SPARQL query đều parse và thực thi thành công.

\section{QA}

Corpus deterministic gồm 20 câu phủ chín intent, yêu cầu đúng intent, nội dung
mong đợi và tối thiểu một evidence record (riêng top-5 yêu cầu đủ năm evidence).
Production path trên Neo4j đạt 20/20. Rubric shortest-path chấp nhận mọi đường hợp
lệ thay vì khóa vào một intermediate; recommendation question kiểm tra contract
và evidence, còn chất lượng ranking được chấm riêng bằng P@10/NDCG@10. Đây vẫn là
smoke corpus do nhóm xây dựng, không phải gold QA được reviewer độc lập xác nhận.

\section{Recommendation}
\reportfigure{recommendation_ablation.pdf}{So sánh kết quả các phương pháp xếp hạng}{fig:recommendation-ablation}

\begin{longtable}{|>{\raggedright\arraybackslash}p{0.22\textwidth}|>{\raggedright\arraybackslash}p{0.22\textwidth}|>{\raggedright\arraybackslash}p{0.22\textwidth}|>{\raggedright\arraybackslash}p{0.22\textwidth}|}
\caption{Tổng hợp recommendation (2)}\label{tab:9-2} \\
\hline
\textbf{Phương pháp} & \textbf{P@10} & \textbf{NDCG@10} & \textbf{Vai trò} \\ \hline
\endfirsthead\hline \textbf{Phương pháp} & \textbf{P@10} & \textbf{NDCG@10} & \textbf{Vai trò} \\ \hline\endhead
Overlap & 0,670 & 0,723 & baseline lịch sử \\ \hline
Weighted Jaccard & 0,640 & 0,699 & baseline lịch sử \\ \hline
Hybrid & 0,590 & 0,657 & baseline lịch sử \\ \hline
IDF-weighted graph & \textbf{0,715} & \textbf{0,754} & production Neo4j \\ \hline
\end{longtable}

IDF đạt kết quả tốt nhất trên 20 case silver và mọi recommendation có explanation.
Tuy nhiên, corpus nhỏ và relevance rubric dựa trên metadata graph; kết quả không
đồng nghĩa người dùng thực sẽ ưa thích recommendation.

\section{Hiệu năng}
\reportfigure{query_latency.pdf}{Phân bố độ trễ truy vấn Neo4j}{fig:query-latency}

Benchmark Neo4j 5.26.28 chạy trên 2.000 Movie, một warm-up và 100 lần/câu. Median
theo intent nằm trong 2,67–188,74 ms; p95 nằm trong 5,08–211,24 ms. Similar-movie
query chậm nhất do candidate traversal và aggregation feature; genre/rating nhanh
nhất. Kết quả mô tả workload và máy đo, chưa chứng minh scalability.

SQLite baseline dùng cùng normalized snapshot cho bốn query tương đương. So sánh
chỉ hợp lệ nếu chạy cùng máy, snapshot, warm-up và iterations. SQLite là baseline
kiểm soát, không đại diện cho mọi RDBMS; mục tiêu là thảo luận trade-off, không
chứng minh graph luôn nhanh hơn relational.

\begin{longtable}{|>{\raggedright\arraybackslash}p{0.18\textwidth}|>{\raggedright\arraybackslash}p{0.18\textwidth}|>{\raggedright\arraybackslash}p{0.18\textwidth}|>{\raggedright\arraybackslash}p{0.18\textwidth}|>{\raggedright\arraybackslash}p{0.18\textwidth}|}
\caption{Tổng hợp hiệu năng (3)}\label{tab:9-3} \\
\hline
\textbf{Intent tương đương} & \textbf{Neo4j median (ms)} & \textbf{Neo4j p95 (ms)} & \textbf{SQLite median (ms)} & \textbf{SQLite p95 (ms)} \\ \hline
\endfirsthead\hline \textbf{Intent tương đương} & \textbf{Neo4j median (ms)} & \textbf{Neo4j p95 (ms)} & \textbf{SQLite median (ms)} & \textbf{SQLite p95 (ms)} \\ \hline\endhead
movies\_by\_director & 10,086 & 16,557 & 1,050 & 1,644 \\ \hline
common\_movies & 38,324 & 48,326 & 12,041 & 13,070 \\ \hline
movies\_by\_genre\_rating & 2,667 & 5,085 & 0,814 & 0,835 \\ \hline
directors\_by\_genre & 4,545 & 6,639 & 2,119 & 2,198 \\ \hline
\end{longtable}

SQLite nhanh hơn trong bốn query kiểm soát này. Kết quả là bằng chứng chống lại
tuyên bố đơn giản “graph luôn nhanh hơn SQL”; lợi ích chính của Neo4j trong đề tài
là mô hình traversal/evidence và một execution surface thống nhất, không phải ưu
thế latency phổ quát.

\section{Threats to validity}

\begin{itemize}
\item Popular-movie sampling tạo selection bias.
\item Top-20 cast bỏ diễn viên phụ.
\item Nguồn TMDB thay đổi theo thời gian; cache chỉ cố định một snapshot.
\item IMDb chỉ enrich Movie, chưa link Person.
\item Silver labels được sinh từ source/rule và chưa độc lập.
\item QA corpus deterministic 20 câu chưa có reviewer độc lập và chưa phủ paraphrase rộng.
\item Recommendation không có user interaction, nên không đánh giá personalization.
\item Benchmark chỉ một quy mô Neo4j và phụ thuộc phần cứng/cache.
\end{itemize}

\section{Trả lời câu hỏi nghiên cứu}

\textbf{RQ1:} Stable source ID và exact IMDb join tích hợp hai nguồn mà không ghi đè
rating; checksum/manifest cung cấp provenance ở mức snapshot.
\textbf{RQ2:} Mười competency question được biểu diễn bằng pattern trực tiếp, multi-hop,
aggregation và shortest path; evidence là node/relationship thực thi.
\textbf{RQ3:} Cypher rule tạo co-star có supporting movie; semantic profile tạo inverse
và type entailment, đồng thời validator phát hiện conflict.
\textbf{RQ4:} IDF giảm ảnh hưởng feature phổ biến, đạt P@10/NDCG@10 cao hơn ba baseline
trên corpus silver và giữ explanation coverage.
\textbf{RQ5:} Hệ thống đạt mục tiêu demo/reproducibility ở 2.000 phim nhưng chưa chứng
minh gold-label accuracy, personalization hoặc multi-scale scalability.

\section{Protocol tái chạy thực nghiệm}

\begin{enumerate}
\item Cố định raw snapshot và ghi checksum.
\item Transform/import authoritative, chạy validation và lưu counts.
\item Chạy reasoning trước evaluation phụ thuộc derived edges.
\item Chạy corpus bằng đúng backend được khai báo.
\item Benchmark một warm-up, 100 iterations/query; không chạy workload khác song song.
\item Lưu Python/Neo4j/platform/movie count trong metadata.
\item Sinh bảng/biểu đồ từ JSON/CSV artifact, không nhập số bằng tay.
\item Kiểm tra report/slide dùng cùng artifact version.
\end{enumerate}

Với relational baseline, snapshot và máy phải giống Neo4j run. Nếu không đáp ứng,
kết quả chỉ được báo riêng, không dùng để xếp hạng engine.

\section{Kiểm thử và release gate}

Lần chạy cuối của \texttt{make test} đạt 30/30 test. Bộ test gồm unit test cho cleaning,
entity resolution, query planning/compiler, RDF export, semantic materialization,
SPARQL catalog, CRUD, relational benchmark và evidence summary. Integration test
dùng Neo4j riêng ở Bolt 7688, thực hiện import hai lần để xác nhận idempotency,
sau đó kiểm tra QA, recommendation và vòng đời CRUD trước khi dọn graph test.

Docker healthcheck gọi \texttt{cypher-shell RETURN 1}, nên trạng thái healthy chỉ được
công bố sau khi Bolt thực sự nhận query; điều này loại race condition trong đó
container đã chạy nhưng driver chưa handshake được. Release gate còn chạy
\texttt{compileall}, dependency check và xác minh checksum của toàn bộ tài liệu nguồn.

\section{Phân tích lỗi và hướng khắc phục}

\subsection{QA}

Rubric cũ khóa shortest path vào một intermediate và đã được thay bằng yêu cầu
đúng intent, có đường/evidence hợp lệ. Recommendation question chỉ kiểm tra API
contract; relevance dùng corpus P@10/NDCG@10 riêng. Bước tiếp theo là bổ sung
paraphrase, câu mơ hồ, expected abstention và review độc lập thay vì tăng điểm
bằng các câu gần template.

\subsection{Entity resolution}

Corpus exact-ID dễ hơn dữ liệu mention thực tế. Cần bổ sung hard negative cùng tên,
alias, Unicode, tên đảo, thiếu ID và candidate tie. Từng error phải phân loại false
merge hoặc missed match vì false merge gây ô nhiễm graph nghiêm trọng hơn.

\subsection{Recommendation}

IDF có thể ưu tiên feature quá hiếm do metadata noise. Có thể đặt minimum document
frequency hoặc review contribution distribution. Tuy nhiên thay đổi trọng số phải
đánh giá trên corpus tách biệt, không điều chỉnh trực tiếp theo 20 case rồi báo
trên chính các case đó.

\subsection{Performance}

Similar-movie chậm nhất gợi ý cần inspect query plan, cardinality và index; nhưng
index không giúp mọi traversal. Cần đo cold/warm cache riêng, nhiều scale thật và
ghi concurrency trước khi tối ưu.
